\begin{textsample}{POS Dim 2 – unprofiled\_gpt – Score 7.00 – unprofiled\_gpt/t000532\_gpt.txt}  \label{ex:f2_pos_028}
I think if you actually gave children really big, knotty historical subjects, you ’d get \textbf{far} \textbf{more} out of them.
Take something like the suffragettes or the world wars – they ’re \textbf{huge}, messy stories with politics, strategy, personalities, social change, all tangled together.
There are loads of \textbf{different} hooks you can get curious about.
You can come at them from \textbf{different} angles, ask why people did what they did, how decisions were made, what might have happened if something had gone differently.
Whereas what we had at school felt like everything was sliced up into these very narrow, self‑contained topics.
You ’d do this tiny little bit of history, in isolation, and it never really went anywhere or connected to anything bigger.
It didn't make me wonder about things, it didn't make me ask questions.
So I just ended up copying stuff off the board, writing down what we were told, instead of actually thinking about it.
I ’m not good at memorisation anyway, and no one ever actually taught us how to memorise, how to work with information so it sticks.
So the whole thing \textbf{became} an exercise in short‑term cramming and regurgitating, \textbf{rather} \textbf{than} learning to think historically.
Which feels like a \textbf{real} shame, given how much my education cost.
It ought to have been better at sparking curiosity \textbf{than} that.

% matched lemmas: become, different, far, huge, more, rather, real, than
\end{textsample}
